\chapter{1903: Svante Arrhenius}
\label{chap:chemistry1903}

The Nobel Prize in Chemistry for the year 1903 was awarded to Svante Arrhenius.

\textbf{Motivation:}
in recognition of the extraordinary services he has rendered to the advancement of chemistry by his electrolytic theory of dissociation

\section{Svante Arrhenius}

\subsection*{Early Life and Education}

Svante Arrhenius was born on 1859-02-19 in Vik, Sweden.
Svante August Arrhenius was a Swedish scientist. Originally a physicist, but often referred to as a chemist, Arrhenius was one of the founders of the science of physical chemistry. In 1903, he received the Nobel Prize in Chemistry, becoming the first Swedish Nobel laureate. In 1905, he became the director of the Nobel Institute, where he remained until his death.

Arrhenius showed an early aptitude for mathematics and science, and he entered the University of Uppsala in 1876. Finding the instruction there conservative and uninspired, he moved in 1881 to the Physical Institute of the Swedish Academy of Sciences in Stockholm, where he worked under the physicist Edvard Olof Edlund. It was there that he carried out the experiments on the conductivity of electrolyte solutions that formed the basis of his doctoral dissertation, submitted to Uppsala in 1884.

\subsection*{Scientific Context}

The behavior of salts, acids, and bases in aqueous solution was one of the central puzzles of chemistry in the nineteenth century. Michael Faraday had shown that solutions of salts conduct electricity, and it was known that the conducting power varied with the substance and with the concentration, but no satisfactory theory explained these facts. At the same time, the study of solutions was being put on a new footing by the work of van 't Hoff, who had shown that dissolved substances exert an osmotic pressure analogous to the pressure of a gas and that the effect depends on the number of dissolved particles. It was into this setting that Arrhenius introduced the bold and controversial idea that salts in solution split into electrically charged fragments, which he called ions.

\subsection*{Career and Major Research}

Arrhenius was professor of physics at Stockholm University and later director of the Nobel Institute for Physical Chemistry. In his 1884 doctoral thesis he proposed the theory of electrolytic dissociation, holding that salts in solution split into charged ions. In 1889 he formulated the equation relating reaction rate to temperature that bears his name. His ideas became central to physical chemistry and electrochemistry.

The theory of electrolytic dissociation met with a hostile reception at first, particularly at the University of Uppsala, where the examining committee regarded the thesis as interesting but unproven. The turning point came in the following years, when van 't Hoff recognized that Arrhenius's theory explained the anomalous osmotic behavior of electrolyte solutions, and Wilhelm Ostwald found that the theory accounted for his measurements of the conductivity of organic acids. With these two powerful allies, Arrhenius was able to develop and defend his ideas, and by the end of the century the theory of electrolytic dissociation was generally accepted.

Arrhenius's interests extended far beyond electrochemistry. In 1889 he published his famous equation relating the rate of a chemical reaction to the temperature, showing that the rate increases exponentially with temperature and introducing the concept of an activation energy, the barrier that must be overcome before a reaction can occur. This relationship, still known as the Arrhenius equation, became one of the cornerstones of chemical kinetics.

\subsection*{Nobel Prize-winning Work}

The work for which Svante Arrhenius was awarded the Nobel Prize in Chemistry in 1903 was described by the Nobel Committee as: "in recognition of the extraordinary services he has rendered to the advancement of chemistry by his electrolytic theory of dissociation".

This research fundamentally advanced the field by introducing the theory of electrolytic dissociation, which explained the conductivity of electrolytes and revolutionized the understanding of acids, bases, and salts in solution.

The essence of the theory was that when a salt such as sodium chloride dissolves in water, it separates into sodium ions and chloride ions, each carrying an electric charge, and that it is the motion of these ions under an applied electric field that carries the current. Arrhenius showed that the electrical conductivity of a solution could be used, with suitable assumptions, to estimate the extent of dissociation for different substances and concentrations. The theory accounted for the osmotic behavior of salts, for the heat changes accompanying neutralization, and for the way in which acids and bases catalyze chemical reactions, and it gave chemists for the first time a coherent picture of what happens when a substance dissolves.

\subsection*{The Work in Detail}

Arrhenius measured the electrical conductivity of a wide range of solutions at different concentrations and showed that the molar or equivalent conductivity of an electrolyte increases with dilution, tending toward a limiting value at very great dilution. He interpreted this by supposing that the fraction of the dissolved substance present as ions increases with dilution, until at infinite dilution the substance is completely dissociated. From such measurements he could estimate, for each concentration, the degree of dissociation and hence the number of particles in the solution, values that agreed with the osmotic pressures measured by van 't Hoff's methods.

The theory also explained a striking regularity observed in the study of acids. It had long been known that the catalytic power of different acids in promoting reactions such as the hydrolysis of sugar varies enormously. Arrhenius and related work by Ostwald connected acid catalysis in dilute solution with the concentration of hydrogen ions, establishing that the hydrogen ion is the active agent in acid catalysis. This result linked the electrical behavior of solutions to their chemical reactivity and gave the theory of dissociation a direct application to the problems of reaction kinetics.

\subsection*{Reception}

The opposition to Arrhenius's theory came largely from chemists who could not accept that a salt in solution could exist as separate charged fragments, since no such separation had ever been demonstrated in the solid state. The controversy was resolved by the combined weight of the conductivity measurements, the agreement with van 't Hoff's osmotic results, and Ostwald's support from the field of acid catalysis. By the early years of the twentieth century the theory was widely accepted and had become a basis of the modern understanding of acids, bases, and solutions. The Nobel Prize of 1903 marked the full rehabilitation of a theory that had received only the lowest passing grade in Arrhenius's doctoral examination of 1884.

\subsection*{Significance and Impact}

The impact of Svante Arrhenius's work extends far beyond the specific achievement recognized by the Nobel Prize.
The theory of electrolytic dissociation transformed the understanding of chemical reactions in solution and provided the framework for modern acid-base theory. Svante Arrhenius's work also pioneered the application of physical principles to chemical problems.

Arrhenius was one of the most wide-ranging scientists of his generation. He applied physical reasoning to an extraordinary variety of problems, from the chemistry of the atmosphere and the oceans to the origin of life and the possibility of life elsewhere in the universe. His studies of the greenhouse effect of carbon dioxide in the atmosphere, published in the 1890s, led him to estimate that a doubling of atmospheric carbon dioxide could raise global temperatures by several degrees, an early anticipation of modern climate science.

\subsection*{Later Career and Honors}

Svante Arrhenius continued his scientific work until 1927-10-02, passing away in Stockholm, Sweden.
He received numerous honors including membership in major scientific academies and societies.
Svante Arrhenius became director of the Nobel Institute for Physical Chemistry in Stockholm and was a key figure in Swedish scientific life.

In 1905 Arrhenius was appointed director of the newly founded Nobel Institute for Physical Chemistry in Stockholm, a position he held until his death. From this post he exercised great influence on the development of Swedish science and participated in the work of the Nobel committees, and he traveled widely, lecturing to scientific audiences throughout Europe and North America. His later researches embraced cosmic physics, the theory of the origin of the solar system, and immunology, and he remained one of the best-known scientific figures in the world until his death in 1927.

\subsection*{Legacy}

The legacy of Svante Arrhenius endures through the lasting impact of his scientific contributions.
The Arrhenius equation describing temperature dependence of reaction rates is ubiquitous in chemical kinetics. Svante Arrhenius's early work on the greenhouse effect anticipated modern climate science.

The two ideas that bear Arrhenius's name, the theory of electrolytic dissociation and the Arrhenius equation, are taught in every course of chemistry and are in daily use in research and industry. The concept of activation energy, which he introduced, remains the central idea of chemical kinetics, and his equation governs the behavior of reactions in fields as different as biochemistry, atmospheric chemistry, and materials science. His early insight into the climatic role of carbon dioxide has been vindicated by a century of observation and is now central to the understanding of global climate change.

\subsection*{Collaborators and Students}

Arrhenius's most important scientific alliances were with Wilhelm Ostwald and Jacobus van 't Hoff, whose work helped establish the discipline of physical chemistry. He collaborated with Ostwald and van 't Hoff, but was not a co-founder of the journal Zeitschrift für physikalische Chemie, which was founded by Ostwald and van 't Hoff. Arrhenius's ideas influenced generations of chemists, and his pupils and colleagues in Stockholm included many who carried his methods into other branches of science. His combination of theoretical insight and wide-ranging curiosity set a pattern for the physical chemist that has persisted to the present day.

\subsection*{Modern Developments}

The Arrhenius equation remains one of the most widely used relationships in chemistry, governing reaction rates in industrial catalysis, atmospheric chemistry, and materials science. Arrhenius's electrolytic dissociation theory is the basis of modern aqueous solution chemistry, pH theory, and electrochemistry. His broader work also anticipated concerns about carbon dioxide and climate change, linking his legacy to present-day climate science.

The modern theory of electrolytes, refined by Peter Debye and Erich Hückel in the 1920s and developed further since, rests in part on the foundations Arrhenius laid, and many measurements of pH, battery technologies, and electrochemical processes draw on concepts he helped establish. In climatology, the connection between atmospheric carbon dioxide and temperature that Arrhenius explored is now recognized as a central mechanism of the Earth's climate, and his early calculations are cited as an important beginning of quantitative climate science.

\subsection*{Selected Publications}

"Recherches sur la conductibilité galvanique des électrolytes" (1884)
"Über die Dissociation der in Wasser gelösten Stoffe" (1887)
"Über die Reaktionsgeschwindigkeit bei der Inversion von Rohrzucker durch Säuren" (1889)
